\documentclass[10pt,a4paper]{amsart}
\usepackage[margin=19mm]{geometry}
\usepackage{amsmath,amssymb,mathtools}
\usepackage[scr=rsfs,cal=euler]{mathalfa}
\usepackage{xfrac}
\usepackage[unicode,hidelinks,bookmarks=false]{hyperref}
\newcommand{\ie}{i.e.\ }
\newcommand{\cf}{cf.\ }
\newcommand{\resp}{respectively\ }
\newcommand{\Subsection}{\subsection}
\providecommand{\nobreakdash}{\nobreak\hbox{-}}
\setlength{\emergencystretch}{2em}
\multlinegap0pt
\usepackage{bm}
\usepackage{bbm}
\usepackage{dsfont}
\usepackage{xspace}
\usepackage{cancel}
\usepackage{mathtools}
\usepackage{amsthm}
\makeatletter
\@ifundefined{Lemma}{\newtheorem{Lemma}{Lemma}}{}
\makeatother
\def\CC{\mathbb C}
\let\myh\widehat
\let\myt\widetilde
\let\myo\overline
\def\tpi{\widetilde{\pi}}
\let\myt\widetilde
\let\myo\overline
\def\tpi{\widetilde{\pi}}
\title[Existence of a "maximal" domain of meromorphy for an analyti]{Existence of a "maximal" domain of meromorphy for an analytic function outside a polar compact set}
\author{Aleksandr Komlov}
\date{Source: arXiv:2608.26497v1 (CC BY 4.0)}
\begin{document}
\maketitle

\noindent\emph{Source and licence.} Existence of a "maximal" domain of meromorphy for an analytic function outside a polar compact set. Version arXiv:2608.26497v1 (CC BY 4.0), distributed under the Creative Commons Attribution 4.0 International licence, as recorded in the arXiv licence metadata for this exact version and shown on the abstract page https://arxiv.org/abs/2608.26497v1. Full rights notice: licenses/arxiv-2608.26497-CC-BY-4.0.txt.\par\smallskip

\noindent\textit{Editorial scope.} Counted unit: the Lemma whose source label is \texttt{L1} in the article (source0.tex lines 450--453, source label \texttt{L1}), delivered together with the article's own complete original proof of it (source0.tex lines 455--469, one \texttt{proof} environment of 15 lines). No mathematics is written, repaired or re-derived: statement and proof are the article's own text line for line. Required context retained uncounted: none. Every definition, display and label of those lines is retained unchanged. Omitted from this package and not redistributed: 1--449, 454--454, 470--884. Nothing inside a retained excerpt is omitted or altered: every retained body is delivered exactly as the article's lines are written, so no omission receipt of any kind accompanies this package. Citations: the retained excerpts carry no citation command at all, so no bibliography entry of the article is transcribed and no reference is redistributed. Reference adaptation: none was needed and none was made; every reference inside the delivered unit resolves inside it, so no unresolved reference or citation is produced. Printed slips kept literal: no printed word, symbol, hypothesis or number of the counted statement or of its proof was altered. Layout and class: the article's own class is \texttt{article} with packages including fontenc, inputenc, amsmath, amsfonts, amssymb, mathrsfs, amscd, amsthm, comment, graphicx, datetime, hyperref; the wrapper is the lane's pinned \texttt{amsart} editorial wrapper at 10pt on A4, which restates the theorem-like environments the retained text needs and adds the article's own macro definitions that the retained text needs (\texttt{\textbackslash CC}, \texttt{\textbackslash myh}, \texttt{\textbackslash myt}), with extra wrapper packages (bm, bbm, dsfont, xspace, cancel, mathtools). The delivered numbering is the unit's own. Assets: no retained span contains a figure environment, a graphics-inclusion command, a PDF-page inclusion, a table or a TikZ picture; no image file is copied into this package, the package declares no asset dependency and no image file is redistributed with it. Licence: version arXiv:2608.26497v1 of this article is distributed under the Creative Commons Attribution 4.0 International licence, as recorded in the arXiv licence metadata for this exact version and shown on the abstract page https://arxiv.org/abs/2608.26497v1. A full rights notice is delivered at licenses/arxiv-2608.26497-CC-BY-4.0.txt. Characters: every retained excerpt is pure ASCII, so the delivered engine, XeLaTeX with its default Latin Modern text font, reports no missing character.
\begin{Lemma}
\label{L1}
Let $f_\infty$ be an analytic germ at $\infty$. Let $\gamma_1$ and $\gamma_2$ be paths from $\infty$ to some point $z_0\in\CC$ such that there exist analytic continuations $f_{z_0}^1$ and $f_{z_0}^2$ of $f_\infty$ along $\gamma_1$ and $\gamma_2$ to $z_0$ respectively and $f_{z_0}^1\ne f_{z_0}^2$. Then there exist a point $\myt z_0\in\CC$ and Jordan non-intersecting at inner points paths $\myt\gamma_1$ and $\myt\gamma_2$ from $\infty$ to $\myt z_0$ such that there exist analytic continuations $f_{\myt z_0}^1$ and $f_{\myt z_0}^2$ of $f_\infty$ along $\myt\gamma_1$ and $\myt\gamma_2$ to $\myt z_0$ respectively and $f_{\myt z_0}^1\ne f_{\myt z_0}^2$. 
\end{Lemma}

\begin{proof}
It immediately follows from the theorem on analytic continuation along close paths that there exist piecewise linear paths $\gamma^\circ_1$ and $\gamma^\circ_2$ such that
there exist analytic continuations of $f_\infty$ to $z_0$ along paths $\gamma^\circ_1$ and $\gamma^\circ_2$, which give $f_{z_0}^1$ and $f_{z_0}^2$, respectively.
Moreover, it is clear that we can choose $\gamma^\circ_1$ and $\gamma^\circ_2$  such that each of them self-intersects at finite number of points and $\gamma^\circ_1$ intersects $\gamma^\circ_2$ also at finite number of points. For simplicity of notation, further we will denote these $\gamma^\circ_1, \gamma^\circ_2$ by$\gamma_1, \gamma_2$.

Let $\gamma_1: [0,1]\to\CC$ self-intersect. Let $\gamma_1(t_2)$ be the smallest (with respect to parameter $t$ of $\gamma_1(t)$) self-intersection point of $\gamma_1$, i.e. $\gamma_1(t_1) = \gamma_1(t_2)$ for some $t_1<t_2$ and $\gamma_1(s_1) \ne \gamma_1(s_2)$ for all $s_1,s_2<t_2$.
Note that $\gamma_1(t), t\in[t_1,t_2]$, is a closed Jordan path.
Let the analytic continuations of $f_\infty$ along $\gamma_1(t), t\in[0,t_1]$, and along $\gamma_1(t), t\in[0,t_2]$, give two different germs at the point $\gamma_1(t_1)=\gamma_1(t_2)$. Then it is clear that we can choose paths $\myt\gamma_1, \myt\gamma_2$  from $\infty$ to $\gamma_1(t_1)$ close to $\gamma_1(t), t\in[0,t_1]$, and $\gamma_1(t), t\in[0,t_2]$, respectively, such that $\myt\gamma_1, \myt\gamma_2$ intersect each other only at their start and end points and the analytic continuation of $f_\infty$ along $\myt\gamma_1$ coincides with the analytic continuation along $\gamma_1(t), t\in[0,t_1]$, and the analytic continuation of $f_\infty$ along $\myt\gamma_2$ coincides with the analytic continuation along $\gamma_1(t), t\in[0,t_2]$. Thus, in this case $\myt\gamma_1, \myt\gamma_2$ are desired paths.
Let now the analytic continuations of $f_\infty$ along $\gamma_1(t), t\in[0,t_1]$, and along $\gamma_1(t), t\in[0,t_2]$, give the same germ at the point $\gamma_1(t_1)=\gamma_1(t_2)$. Then we can cut the loop $\gamma_1(t), t\in[t_1, t_2]$. Indeed, the path $\gamma_1(t), t\in[0,t_1]\cup [t_2,1]$, has strictly fewer self-intersection points than $\gamma_1(t), t\in[0,1]$, and the analytic continuation along both these paths give the same germ $f_{z_0}^1$. Now we can repeat all this procedure for the path $\gamma_1(t), t\in[0,t_1]\cup [t_2,1]$, an so on. As a result we get either desired paths $\myt\gamma_1, \myt\gamma_2$ or a Jordan path $\breve\gamma_1$ from $\infty$ to $z_0$ such that there exists the analytic continuation of $f_\infty$ along $\breve\gamma_1$ that gives the germ $f_{z_0}^1$.

Now we do the same procedure for $\gamma_2$. As a result we have either desired paths $\myt\gamma_1, \myt\gamma_2$ or two Jordan paths $\breve\gamma_1, \breve\gamma_2$ from $\infty$ to $z_0$ such that $\breve\gamma_1, \breve\gamma_2$ intersect each other at finite numbers of points and there exist the analytic continuations of $f_\infty$ along $\breve\gamma_1$ and $\breve\gamma_2$ that give the germs $f_{z_0}^1$ and $f_{z_0}^2$, respectively.
Let us consider the second situation. Let $\breve\gamma_1: [0,1]\to\myh\CC$ and $\breve\gamma_2: [0,1]\to\myh\CC$. Let $t_1, s_1\in(0,1]$ be such that $\breve\gamma_1(t_1)=\breve\gamma_2(s_1)$ and $\breve\gamma_1(t)\ne\breve\gamma_2(s)$ for all $t<t_1$, $s\in(0,1]$ and $\breve\gamma_1(t_1)\ne\breve\gamma_2(s)$ for all $s<s_1$. 
Note that $\breve\gamma_1(t), t\in[0,t_1]$, and $\breve\gamma_2(t), t\in[0,s_1]$, are two Jordan paths from $\infty$ to $\breve\gamma_1(t_1)=\breve\gamma_2(s_1)$ that don't intersect each other at inner points. Therefore if the analytic continuations of $f_\infty$ along $\breve\gamma_1(t), t\in[0,t_1]$, and $\breve\gamma_2(t), t\in[0,s_1]$, give two different germs, these paths are desired paths $\myt\gamma_1, \myt\gamma_2$. 
Let us suppose that the analytic continuations of $f_\infty$ along $\breve\gamma_1(t), t\in[0,t_1]$, and $\breve\gamma_2(t), t\in[0,s_1]$, give the same germ. Then the analytic continuation of $f_\infty$ along $\breve\gamma_2(t), t\in[0,1]$, coincides with the analytic continuation along the path $\breve\gamma_3$ that is the union of $\breve\gamma_1(t), t\in[0,t_1]$, and $\breve\gamma_2(t), t\in[s_1,1]$. It is clear that we can choose a path $\breve\gamma_4$ from $\infty$ to $z_0$ close to $\breve\gamma_3$ such that the analytic continuations along them coincide, $\breve\gamma_4$ is the union of some path $\breve\gamma$ from $\infty$ to $\breve\gamma_2(s_0)$ for some $s_0>s_1$ and $\breve\gamma_2(t), t\in[s_0,1]$, and the paths $\breve\gamma$ and $\breve\gamma_1$ don't intersect each other. So, the analytic continuation of $f_\infty$ along $\breve\gamma_4$ also as along $\breve\gamma_2$ gives the germ $f_{z_0}^2$, but the number of intersection points of $\breve\gamma_4$ and $\breve\gamma_1$ is strictly fewer than of $\breve\gamma_2$ and $\breve\gamma_1$. Repeating this procedure several times we get the desired paths $\myt\gamma_1, \myt\gamma_2$.
\end{proof}

\end{document}
